\subsection{Biến mệnh đề và quy ước ký hiệu}
\label{sec:prop_vars}

Để đưa bài toán về logic mệnh đề, mỗi ô $i \in \mathcal{C}$ được gắn một biến
mệnh đề $x_i$ với ngữ nghĩa
\begin{equation}
  x_i = \top \iff \text{nhân vật tại ô } i \text{ là tội phạm}.
  \label{eq:prop_var}
\end{equation}
Literal $x_i$ khẳng định ô $i$ là tội phạm, literal $\neg x_i$ khẳng định ô $i$
vô tội. Một phép gán đầy đủ cho $\{x_i\}_{i\in\mathcal{C}}$ tương ứng một-một
với một cấu hình tội trạng của bàn chơi.

\paragraph{Đánh số biến cho bộ giải SAT.}
Bộ mã hóa \texttt{CNFEncoder} biểu diễn literal theo quy ước DIMACS: biến là
số nguyên dương, literal âm là phủ định. Các định danh ô được sắp tăng dần rồi
đánh số từ $1$:
\begin{equation}
  v(i) = \bigl|\{\, j \in \mathcal{C} : j < i \,\}\bigr| + 1 .
  \label{eq:dimacs_var}
\end{equation}
Vì định danh ô liên tiếp từ $0$ nên $v(i) = i + 1$ và
$v(\mathcal{C}) = \{1, \ldots, |\mathcal{C}|\}$. Literal $+v(i)$ biểu diễn
$x_i$, literal $-v(i)$ biểu diễn $\neg x_i$.

\paragraph{Biến phụ trợ.}
Hai dạng ràng buộc cần biến không gắn với ô nào. Bộ mã hóa cấp phát chúng từ
bộ đếm \texttt{next\_variable}, khởi tạo bằng $|\mathcal{C}|+1$ và tăng sau mỗi
lần cấp phát, nên biến phụ trợ không bao giờ trùng biến ô:
\begin{itemize}
  \item \texttt{ParityConstraint} trên $m$ literal dùng $m-1$ biến $p_k$ lưu
    tổng XOR tiền tố;
  \item \texttt{ConnectedConstraint} trên $m$ ô dùng $m^2$ biến khả đạt
    $r_k^{(t)}$, với $m$ biến cho mỗi bước $t = 0, \ldots, m-1$.
\end{itemize}
Định nghĩa chi tiết của hai nhóm biến này nằm ở Mục~\ref{sec:clue_semantics}.
Bảng~\ref{tab:ky_hieu} tổng hợp các ký hiệu dùng trong chương.

\begin{table}[htbp]
  \centering
  \caption{Ký hiệu dùng trong mô hình hóa bài toán}
  \label{tab:ky_hieu}
  \begin{tabularx}{\textwidth}{l>{\raggedright\arraybackslash}X}
    \toprule
    \textbf{Ký hiệu} & \textbf{Ý nghĩa} \\
    \midrule
    $H, W$                 & Số hàng và số cột của bàn chơi \\
    $\mathcal{C}$          & Tập định danh ô, $|\mathcal{C}| = HW$ \\
    $x_i$                  & Biến mệnh đề của ô $i$; $x_i=\top$ khi ô $i$ là tội phạm \\
    $v(i)$                 & Chỉ số biến DIMACS của ô $i$ \\
    $\sigma^\star$         & Lời giải khai báo của bản đồ \\
    $\mathcal{K}$          & Tập toàn bộ lời khai (clue) của bản đồ \\
    $\mathrm{src}(\gamma)$ & Ô sở hữu lời khai $\gamma \in \mathcal{K}$ \\
    $\phi_\gamma$          & Tập mệnh đề CNF mã hóa ràng buộc của $\gamma$ \\
    $\mathcal{R}_t$        & Tập lời khai đã lật mở sau $t$ phán quyết được chấp nhận \\
    $\mathcal{P}_t$        & Tập phán quyết đã được chứng minh sau $t$ phán quyết \\
    $\mathit{KB}_t$        & Cơ sở tri thức công khai tương ứng với $(\mathcal{R}_t, \mathcal{P}_t)$ \\
    \bottomrule
  \end{tabularx}
\end{table}
